\section{Conclusion（結論）}

本研究では、ホテルサンバレー那須における地域住民参加型クリスマス聖歌企画を通じ、

\begin{itemize}
\item 地域住民
\item 観光客
\item ホテル
\end{itemize}

の三者が共同で観光体験を形成する「相互創生構造」を分析した。

その結果、観光は単なる消費活動ではなく、「共に時間を過ごし、共に空間を作る体験」として再構成可能であることが示された。

また本研究では、

\begin{itemize}
\item 教会空間
\item オルガン
\item 地域住民
\item 高齢者
\item 音楽
\item 文化
\item 観光客
\end{itemize}

など、地域に既に存在していた素材が接続されることで、新たな観光価値が形成されていた。

これは、地域社会に必要なのが常に新しいものを作り続けることではなく、

\begin{quote}
「既に存在している価値を接続し直すこと」
\end{quote}

である可能性を示している。

さらに、その接続を支えていたのは、

\begin{itemize}
\item 歌詞カード
\item 曲解説
\item 感謝メッセージ
\item 参加導線
\end{itemize}

などによる「情報設計」であった。

本研究で観察された構造は、単なるイベント設計ではなく、人・文化・空間・感情を接続する社会システムとして機能していた。

そして、そのような接続構造を社会へ丁寧に実装していくことこそ、これからの社会インフラとしてのITの役割である可能性がある。

最後に、本研究で最も重要だったのは、地域住民・ホテル関係者・観光客が、それぞれ異なる立場でありながら、「同じ時間を共に過ごした」という感覚である。

共に歌い、共に耳を澄まし、クリスマスの時間を分かち合うことで、

\begin{itemize}
\item 訪れる人
\item 迎える人
\item 暮らす人
\end{itemize}

の境界は静かに薄れ、ホテルサンバレー那須という場所そのものが、一つの地域文化空間として機能していた。

そして、この小さな地方の中でこそ、人・文化・空間・感情といった既に存在している情報を丁寧に整理しやすく、その情報の組み合わせを変えることで生まれる接続性が見え始めていた。

そのような接続構造は、同じ課題を抱える地域社会における新たな地域モデルとなり、さらにその成功が積み重なっていくことで、やがて都市社会における共同体再構築のテンプレートとなっていくであろう。
